%============================================================
% KATA PENGANTAR
%   Diakhiri kota, tanggal, dan kata "Penulis" rata kiri di
%   bawahnya (tanpa nama terang & tanda tangan) — sesuai template
%   asli. Spasi 2, paragraf menjorok 1,27cm.
%   Teks di bawah disalin verbatim dari template asli
%   (reference/Template_Skripsi.docx) — ganti dengan kata
%   pengantar Anda.
%============================================================
\chapter*{KATA PENGANTAR}
\TocHalaman{KATA PENGANTAR}
\setstretch{1.9}

Nama harus ditulis secara lengkap termasuk gelar akademik dan
harus dihindari ucapan terima kasih kepada pihak yang tidak
terkait. Bahasa yang digunakan harus mengikuti kaidah bahasa
Indonesia yang baku.

\noindent
Bagian ini tidak perlu dituliskan hal-hal yang bersifat ilmiah.
Kata Pengantar diakhiri dengan mencantumkan kota dan tanggal
penulisan diikuti di bawahnya dengan
\NamaTTD{kata ``Penulis'' tanpa perlu menyebutkan nama dan tanda
tangan}.

% Penutup: kota + tanggal, diikuti kata "Penulis" di bawahnya —
% keduanya RATA KIRI (sesuai template asli, tanpa nama terang dan
% tanpa tanda tangan).
\vspace{0.8em}
\noindent
\TanggalKataPengantar

\vspace{0.6em}
\noindent
Penulis

\clearpage
